\documentclass[11pt,a4paper]{article}
\usepackage[utf8]{inputenc}
\usepackage[T1]{fontenc}
\usepackage{amsmath,amssymb,amsthm}
\usepackage{geometry}
\geometry{margin=2.4cm}

\newtheorem{definition}{Definition}
\newtheorem{theorem}{Theorem}
\newtheorem{remark}{Remark}

\title{The Critical Line as a Stability Boundary Interpretation: An OVRP Audit Across Symmetry, Oscillation, and Open Foundations}
\author{Mariano Panzano Caball\'e \\ \small Oasis Sovereign Architecture Research, Barcelona}
\date{November 25, 2026}

\begin{document}
\maketitle

\begin{abstract}
Following the integrated compendium (\texttt{DOI: 10.5281/zenodo.24539182}) under the \emph{Oasis Verifiable Research Pipeline} (OVRP, \texttt{DOI: 10.5281/zenodo.24443531}), this note (\texttt{Paper~25}) explores an interpretive physical reading of the critical line $\operatorname{Re}(s) = \tfrac{1}{2}$ as a conceptual boundary between oscillatory phase and damping/stability balance. Rather than claiming a proof of the Riemann Hypothesis or a complete theory of everything, we pair a zero-\texttt{sorry} \texttt{Lean~4.36.0} specification (\texttt{Paper25\_CriticalLine\_Spec.lean}) on finite midpoint symmetry and bounded balance with a deterministic \texttt{C99} damping/resonance probe sealed via \texttt{SHA-256} and Bitcoin OpenTimestamps.
\end{abstract}

\section{Interpretive Boundary and Finite Audit}
Mathematically, the Riemann zeta function $\zeta(s)$ on $s = \sigma + it$ connects to prime distribution, while the Riemann Hypothesis---whether all non-trivial zeros lie on $\sigma = \tfrac{1}{2}$---remains open. In an interpretive, divulgative register, viewing $\sigma = \tfrac{1}{2}$ as a ``stability boundary'' illustrates balance between competing regimes, provided finite checks and open mathematical/physical problems are kept strictly distinct.

\begin{remark}[Explicit Open Boundaries]
An interpretive stability picture and finite \texttt{Lean~4.36.0} / \texttt{C99} checks do not prove the Riemann Hypothesis, continuum Yang--Mills confinement, or a unified theory of fundamental interactions.
\end{remark}

\end{document}
